\documentclass[11pt,a4paper,DIV=10,parskip=half,BCOR=0mm]{scrartcl}
\usepackage[utf8]{inputenc}
\usepackage[T1]{fontenc}
\usepackage[ngerman]{babel}
\usepackage{scrpage2} % pagestyle
\usepackage{geometry} % Seitenränder
\usepackage{amsmath}
\usepackage{amsfonts}
\usepackage{amssymb}
\usepackage{amsthm}
\usepackage{multicol}
\usepackage{graphicx}
\usepackage{float}
\usepackage{marginnote} % Randnotizen
\usepackage{tabularx} % Tabelle mit spaltentyp 'X' für feste breite

% Seitenränder einstellen: marginparwidth wichtig für Randnotizen!
\geometry{a4paper, top=25mm, left=35mm, right=50mm, bottom=30mm,
	headsep=10mm, footskip=12mm, marginparwidth=45mm, marginparsep=3mm}

% Fußzeile
\pagestyle{scrheadings}
\ifoot{ \hrule {\textcopyright} \ C. van Heteren-Frese}

\begin{document}
\pagenumbering{gobble}
%
% === HEADER ====
\setlength{\tabcolsep}{0mm} % kein Innenrand bei Spaleten
\begin{tabularx}{\linewidth}{lXr}
{\Large\textsf{\textbf{Gruppe C:} Lobby \textit{Pro Kernenergie}}} & & \includegraphics[scale=0.04]{images/atom}\\
\hline
\end{tabularx}
% === HEADER END ===
%
\subsection*{Hintergrund}
Du bist Mitglied der Atomlobby \marginnote{\footnotesize \textbf{Lobby:} Gruppe die ein gemeinsames Interesse verfolgt} und vertrittst einen Zusammenschluss
verschiedener Firmen die Atomkraftwerke betreiben. Ihr verdient damit
eine Menge Geld und wollt das natürlich auch in Zukunft tun.

Außerdem setzt Du dich mit deiner Organisation für den Datenschutz und
Verbraucherrechte ein, um in der Öffentlichkeit gut darzustehen.
\subsection*{Eure Position}
\begin{enumerate}
	\item[•]Ihr seid \textbf{gegen} das Smart Grid
\end{enumerate}
\subsection*{Eure Gründe}
\begin{enumerate}
	\item[•]Ihr behauptet, dass erneuerbare Energien zu unsicher sind und das der Strom knapp wird, wenn das Wetter nicht mitspielt.
	\item[•]Wenn mit Hilfe des Smart Grids immer mehr erneuerbare Energien
genutzt werden können, wird weniger Atomstrom genutzt!
	\item[•] Wenn weniger Atomstrom genutzt wird, verdient ihr weniger Geld.
Dagegen müsst ihr was unternehmen!
\end{enumerate}
\end{document}